\documentclass[11pt]{article}
\usepackage{amsmath,amssymb,amsthm,hyperref}
\newtheorem{theorem}{Theorem}
\newtheorem{lemma}{Lemma}
\newcommand{\N}{\mathbb N}
\newcommand{\Z}{\mathbb Z}
\title{Periodicity of errors after reducing rational tails}
\author{Research notes on Erd\H{o}s 243; no priority claim}
\date{6 November 2026}
\begin{document}
\maketitle

\noindent\textbf{Status.} This is a restricted theorem, not a solution
of Erd\H{o}s 243. It rules out eventual periodicity and two particular
bounded alphabets of signed errors in lowest terms in a counterexample.
It allows infinitely many common
factors to be cancelled when reducing successive tails. It does not
assert that bounded or arbitrary errors are periodic.

The rational-tail framework is from
\href{https://math.colgate.edu/~integers/aa28/aa28.pdf}{Koizumi (2026)}.
The use of large prime powers in a reduced denominator, and their
protection below a numerator threshold, is developed in
\href{https://github.com/wcook04/plectis-erdos/blob/main/paper/243/erdos243-reciprocal-tail-reasoning-surface.tex}{Cook's companion paper}.
We prove explicitly the variants needed here, including exclusion of
any fixed finite set of primes from the supply lemma. No priority is
claimed.

\paragraph{Reduced states.}
Let $a_n$ be strictly increasing positive integers with
$a_{n+1}/a_n^2\to1$ and rational reciprocal sum. Write the positive tail
in lowest terms as
\[
 R_n=\sum_{j\ge n}\frac1{a_j}=\frac{u_n}{v_n},\qquad
 \gcd(u_n,v_n)=1.
\]
Define integers
\[
 e_n=v_n-(a_n-1)u_n,\quad
 w_n=a_nu_n-v_n=u_n-e_n>0,\quad
 h_n=\gcd(w_n,a_nv_n).
\]
Then
\begin{equation}\label{eq:reduced}
 h_nu_{n+1}=u_n-e_n,\qquad h_nv_{n+1}=a_nv_n.
\end{equation}
The usual tail expansion gives $v_n/u_n=a_n-1+o(1)$, so
$e_n/u_n\to0$. Also, for some $c>0$,
\begin{equation}\label{eq:growth}
 \log a_n=c2^n+o(1),\qquad \log u_n=o(n).
\end{equation}
Here the first assertion follows by summing
$\log a_{n+1}-2\log a_n=o(1)$ with geometric weights; a zero limiting
coefficient would imply $\log a_n\to0$, contrary to strict increase.
For the second, if the full sum has denominator $q$, the unreduced
integer numerator $C_n=q(\prod_{j<n}a_j)R_n$ has
$C_{n+1}/C_n=a_n-R_n^{-1}\to1$. Hence $\log C_n=o(n)$ and
$u_n\le C_n$ gives the assertion.

An error zero at a sufficiently late index is absorbing. Indeed,
$e_n=0$ implies $u_n=1$, $v_n=a_n-1$, and the next reduced tail is
$1/[a_n(a_n-1)]$. Thus $u_{n+1}=1$; since $e_{n+1}/u_{n+1}\to0$
and the next errors are integers, on a tail where all relative errors
have absolute value below $1$ this forces $e_{n+1}=0$, and induction
applies. It also gives $a_{n+1}=a_n^2-a_n+1$. Therefore, in a
counterexample, every sufficiently late $e_n$ is nonzero and
\begin{equation}\label{eq:udiverges}
 u_n\longrightarrow\infty
\end{equation}
by $1/u_n\le|e_n|/u_n\to0$.

\begin{lemma}[Prime powers outside a finite set]\label{lem:supply}
For any finite set $S$ of primes and any fixed positive integer $K$,
all sufficiently late $v_n$ have a prime-power divisor $Q=p^k$ with
$p\notin S$ and $Q>u_n+K$.
\end{lemma}
\begin{proof}
The reduced denominator $v_n$ divides
$L_n=\operatorname{lcm}(q,a_1,\ldots,a_{n-1})$. Put $s=|S|$.
For each prime in $S$, its maximal exponent in this lcm is attained
in $q$ or one earlier multiplier. After merging repeated indices, the
$S$ part of $v_n$ is at most $q$ times a product of at most $s$
distinct earlier multipliers, and hence at most
$q\prod_{j=1}^s a_{n-j}$. By~\eqref{eq:growth} its logarithm is at most
$c(1-2^{-s})2^n+O(1)$. On the other hand,
$\log v_n=\log u_n+\log a_n+o(1)\ge c2^n+o(1)$.
Consequently the part $v_n^{(S)}$ supported outside $S$ satisfies
\[
 \log v_n^{(S)}\ge c_S2^n
\]
for some $c_S>0$ and all sufficiently large $n$.
If all its prime-power factors were at most $B_n=u_n+K$, it would
divide $\operatorname{lcm}(1,\ldots,B_n)$, so
$\log v_n^{(S)}\le B_n\log B_n$. But
$\log B_n=o(n)$ by~\eqref{eq:growth}, so
$B_n\log B_n=e^{o(n)}=o(2^n)$, a contradiction.
\end{proof}

\begin{lemma}[Protection of a prime power]\label{lem:protect}
In a reduced step~\eqref{eq:reduced}, if $p^k\mid v_n$ and
$0<w_n<p^{k+1}$, then $p^k\mid v_{n+1}$.
\end{lemma}
\begin{proof}
Write $b=\nu_p(v_n)\ge k$ and $d=\nu_p(a_n)$. Coprimality makes
$u_n$ a unit modulo $p$. If $d<b$ then $\nu_p(w_n)=d$ and
$\nu_p(v_{n+1})=b$; if $d>b$ then
$\nu_p(w_n)=b$ and $\nu_p(v_{n+1})=d$. If $d=b$, then
$\nu_p(w_n)\ge b$, whereas $w_n<p^{k+1}$ makes
$\nu_p(w_n)\le k$. Thus $b=d=k$ and $\nu_p(w_n)=k$, again giving
$\nu_p(v_{n+1})=k$. These cases exhaust the possibilities.
\end{proof}

\begin{theorem}[Periodic reduced errors force a Sylvester tail]
Under the stated hypotheses, if $(e_n)$ is eventually periodic, then
\[
 a_{n+1}=a_n^2-a_n+1
\]
for every sufficiently large $n$.
\end{theorem}
\begin{proof}
Assume there is no Sylvester tail and shift to a tail where the errors
are periodic with period $L\ge1$, nonzero, and $|e_n|\le M$ for a
fixed $M\ge1$. Then~\eqref{eq:udiverges} holds. Write
$T=\sum_{i=0}^{L-1}e_i$ for the signed period sum.

First suppose $T=0$. During any consecutive run of steps with $h_n=1$,
the numerator changes by minus the partial sums of the periodic errors.
Thus it never rises more than $LM$ above its value at the start of the
run. At a step with $h_n\ge2$,
$u_{n+1}\le(u_n+M)/2$. If $B$ is the value just after one such
step, the value just after the next is at most
$(B+(L+1)M)/2$. These bounds also hold from the initial value of the
tail. Therefore all run-start values are bounded by
$\max(u_N,(L+1)M)$, and all values are bounded by this plus $LM$.
The same conclusion is immediate if there are only finitely many
cancellations. This contradicts $u_n\to\infty$.

Now suppose $T\ne0$. Let $S$ be the finite set of prime divisors of
$|T|$, with $S$ empty if $|T|=1$. Choose a sufficiently late index $N$
and, by Lemma~\ref{lem:supply}, a divisor $Q=p^k$ of $v_N$ such that
\[
 p\nmid T,\qquad Q>u_N+2ML+1.
\]
We show $Q$ can never be lost. Work up to a putative first step at
which $p^k$ ceases to divide the next reduced denominator. Before it,
$p\mid v_n$ and therefore $p\nmid u_n$.

There cannot be $Lp$ consecutive steps with $h_n=1$ in this interval.
In such a run, taking every $L$th state gives
$u_{t+jL}=u_t-jT$ modulo $p$. Since $p\nmid T$, one of the first
$p$ such states is zero modulo $p$, contrary to coprimality.

Put $H=Lp$. Between cancellations there are at most $H-1$ steps,
each increasing $u$ by at most $M$. If a run starts at $B$, its
next post-cancellation value is at most $(B+HM)/2$. Inductively every
run-start value, up to and including the putative loss step, is at most
$\max(u_N,HM)$. At every intervening step, including that loss step,
\[
 w_n=u_n-e_n\le\max(u_N,HM)+HM
 \le u_N+2MLp<pQ=p^{k+1}.
\]
The final strict inequality follows from the choice of $Q$ and $p\ge2$.
Lemma~\ref{lem:protect} contradicts loss of $p^k$. Thus no such
first loss exists. The same run-length and contraction estimates now
bound $u_n$ for all $n\ge N$, again contradicting~\eqref{eq:udiverges}.
This finishes both cases.
\end{proof}

\begin{theorem}[Fixed magnitude, with arbitrary signs]\label{thm:magnitude}
Under the same rational near-squaring hypotheses, if $|e_n|=m$ for
a fixed $m\ge1$ and every sufficiently large $n$, then a Sylvester
tail follows. Thus a nonterminal error sequence cannot eventually have
a fixed nonzero absolute value, even with an aperiodic sign pattern.
\end{theorem}
\begin{proof}
Suppose a nonterminal orbit has $e_n\in\{-m,m\}$ on a tail. Then
$u_n\to\infty$. Apply Lemma~\ref{lem:supply} with the prime divisors
of $m$ as the excluded set and choose $Q=p^k\mid v_N$ such that
$p\nmid m$ and $Q>u_N+2m+1$.

Work up to a putative first loss of $p^k$. In a consecutive run with
$h_n=1$, the numerator takes steps $\pm m$ on the lattice
$B+m\Z$, starting at its initial value $B$. All vertices are units
modulo $p$, because the denominator is divisible by $p$. There is a
multiple of $p$ on this lattice above $B$, at distance at most $mp$.
A nearest-neighbor walk on the lattice cannot pass it without landing
on it. Thus throughout the run $u_n\le B+m(p-1)$, a bound valid
regardless of the run's length. At the next cancellation,
\[
 u_{n+1}\le\frac{B+mp}{2}.
\]
Consequently all run-start values before first loss are at most
$\max(u_N,mp)$, and at every step, including the putative loss step,
\[
 w_n\le\max(u_N,mp)+mp\le u_N+2mp<pQ.
\]
Lemma~\ref{lem:protect} forbids first loss. The same bounds then hold
forever and contradict $u_n\to\infty$.
\end{proof}

\begin{theorem}[Two magnitudes with odd common factor]\label{thm:twomagnitudes}
Under the same hypotheses, fix an odd positive integer $m$. If
\[
 e_n\in\{-2m,-m,m,2m\}
\]
for every sufficiently large $n$, then a Sylvester tail follows.
No periodicity of the choices or signs is assumed.
\end{theorem}
\begin{proof}
Suppose this is a nonterminal orbit, so $u_n\to\infty$. Exclude the
prime divisors of $2m$ and choose, by Lemma~\ref{lem:supply},
$Q=p^k\mid v_N$ with $p\nmid2m$ and $Q>u_N+6m+1$.

Before a putative first loss, consider a run with $h_n=1$ starting at
$B$. The numerators lie in $B+m\Z$ and take steps of length $m$
or $2m$, in either direction. Moreover
\[
 \gcd(u_n,e_n)=\gcd(u_n,v_n)=1.
\]
In particular, a step of length $2m$ must start at an odd integer.
There is an odd multiple $z$ of $p$ on $B+m\Z$ above $B$, at
distance at most $2mp$: multiples of $p$ on this lattice occur at
spacing $mp$, which is odd, so their parities alternate.

This $z$ is an impassable upper barrier for the run. Landing on it
would contradict $p\nmid u_n$. A step of length $m$ cannot skip a
lattice vertex. The only step of length $2m$ that could pass $z$
without landing is from $z-m$ to $z+m$; since $m$ and $z$ are odd,
it would start at an even integer, contradicting
$\gcd(u_n,e_n)=1$. Hence $u_n<B+2mp$ throughout the run.
At its next cancellation, $w_n\le B+2mp+2m$. Put
$J=2m(p+1)\le3mp$ since $p\ge2$. Post-cancellation values are
therefore at most $(B+J)/2$, and induction bounds every run-start
value by $\max(u_N,J)$. At every step up to and including first loss,
\[
 w_n\le\max(u_N,J)+J\le u_N+2J\le u_N+6mp<pQ.
\]
Protection again prevents first loss, and the resulting uniform bound
on all $u_n$ contradicts $u_n\to\infty$.
\end{proof}

\paragraph{Limit of the result.}
A possible counterexample must therefore have reduced errors that are
not eventually periodic. Its magnitudes cannot eventually be constant,
or confined to $\{m,2m\}$ for a fixed odd $m$.
The proof does not exclude all bounded aperiodic integer error sequences,
nor unbounded errors whose relative size tends to zero. These remaining cases have not been excluded
or constructed here. The full proof or disproof of Erd\H{o}s 243 is
still outstanding.
\end{document}
